\chapter{Resultados --- Zaida (requisitos, frontend y experiencia de usuario)}
\label{cap:resultados-zaida}

Este cap\'itulo integra la informaci\'on real disponible sobre el trabajo de
Zaida: ingenier\'ia de requisitos, frontend Angular, accesibilidad manual
verificable en el c\'odigo, e integraci\'on con el esquema de autenticaci\'on
por cookies. Las mediciones cuantitativas de usabilidad (SUS) y
accesibilidad autom\'atica (Lighthouse) \textbf{no se reportan como
resultados} porque a\'un no se han ejecutado; \'unicamente se describe su
metodolog\'ia (cap\'itulo~\ref{cap:protocolo-experimental}).

\section{Requisitos}

El repositorio incluye una consolidaci\'on completa de ingenier\'ia de
requisitos: \texttt{docs/requisitos/SRS.md} (especificaci\'on de
requisitos de software), \texttt{docs/requisitos/historias/HistoriasUsuario.md},
\texttt{docs/requisitos/casos-de-uso/CasosDeUso.md} y
\texttt{docs/requisitos/CHANGELOG-REQ.md}. La matriz de trazabilidad
derivada (\texttt{docs/trazabilidad/matriz.csv}) vincula cada requisito con
su historia de usuario, caso de uso, m\'odulo de c\'odigo, endpoint y
prueba automatizada (cap\'itulo~\ref{cap:trazabilidad}).

\section{Frontend Angular}

Aplicaci\'on Angular~17.3 con componentes \emph{standalone} (sin
\texttt{NgModule}) y formularios reactivos (\texttt{ReactiveFormsModule}).
M\'odulos verificados directamente en \texttt{frontend/src/app/}:

\begin{itemize}[nosep]
  \item \texttt{core/auth.service.ts} / \texttt{core/auth.guard.ts}:
    integraci\'on con el esquema de cookies del backend
    (\texttt{credentials: include}), sin JWT en \texttt{localStorage}, y
    protecci\'on de rutas que requieren sesi\'on.
  \item \texttt{core/http-error.interceptor.ts} /
    \texttt{core/problem-detail.service.ts}: interceptan las respuestas de
    error y traducen el \texttt{ProblemDetail} del backend a un mensaje
    legible para el usuario.
  \item \texttt{features/login.component.ts}: formulario de login con
    validaci\'on reactiva y manejo de foco accesible.
  \item \texttt{features/mascotas.component.ts} /
    \texttt{features/mascota-api.service.ts}: listado, alta, edici\'on y
    baja de mascotas contra la API real.
\end{itemize}

% EVIDENCIA-ZAIDA-01
% Ruta esperada: figuras/zaida/01-frontend-login.png
\evidencia{figuras/zaida/01-frontend-login.png}%
{Pantalla de inicio de sesi\'on del frontend de BIOPET.}%
{fig:zaida-frontend-login}

% EVIDENCIA-ZAIDA-02
% Ruta esperada: figuras/zaida/02-crud-mascotas.png
\evidencia{figuras/zaida/02-crud-mascotas.png}%
{CRUD de mascotas en el frontend: listado, alta y edici\'on.}%
{fig:zaida-crud-mascotas}

\section{Autenticaci\'on mediante cookies (integraci\'on frontend)}

El frontend no maneja el JWT directamente: env\'ia y recibe las cookies
\texttt{access_token}/\texttt{refresh_token} de forma transparente gracias
a \texttt{credentials: include} en las peticiones HTTP, consistente con el
esquema \texttt{HttpOnly}+\texttt{Secure}+\texttt{SameSite=Strict} descrito
en el cap\'itulo~\ref{cap:resultados-jaime}. Esta integraci\'on qued\'o
reflejada en el historial de cambios del proyecto (rama
\texttt{zaida/frontend-cookies-jwt} y, m\'as recientemente,
\texttt{zaida/problemdetail-409-accesibilidad-login}, que ajust\'o el
manejo de \texttt{ProblemDetail} y la accesibilidad del login tras la
introducci\'on del c\'odigo 409 en el registro).

\section{CRUD accesible de mascotas y accesibilidad manual}

El c\'odigo fuente confirma el uso deliberado de atributos ARIA en los
formularios cr\'iticos: \texttt{login.component.ts} usa
\texttt{[attr.aria-invalid]}, \texttt{[attr.aria-describedby]} vinculado al
mensaje de error de cada campo, y un bloque de error con
\texttt{role="alert"} y \texttt{aria-live="assertive"} para que un lector
de pantalla anuncie el error sin que el usuario deba buscarlo
manualmente, en l\'inea con las pautas de accesibilidad WCAG~\cite{wcag21}
--- el color rojo del CSS es un refuerzo visual adicional, no
el \'unico canal del mensaje. El componente de mascotas
(\texttt{mascotas.component.ts}) usa el mismo tipo de atributos de forma
extensiva en su marcado. Estas verificaciones son de c\'odigo, no de una
herramienta automatizada de accesibilidad.

% EVIDENCIA-ZAIDA-03
% Ruta esperada: figuras/zaida/03-responsive.png
\evidencia{figuras/zaida/03-responsive.png}%
{Vista responsive del frontend en un ancho de pantalla reducido.}%
{fig:zaida-responsive}

\section{Usabilidad (SUS)}

El instrumento System Usability Scale (Brooke~\cite{sus-brooke1996}),
alineado con los conceptos de usabilidad de ISO~9241-11~\cite{iso9241-11},
con un m\'inimo de diez
participantes externos al equipo, est\'a documentado metodol\'ogicamente en
\texttt{docs/etica/ETHICS.md} (secci\'on~iii): consentimiento informado por
participante, tarea de referencia (login, alta, edici\'on, eliminaci\'on
l\'ogica, logout), y resultados agregados vinculados \'unicamente a un
c\'odigo de participante. \textbf{A la fecha de este informe no existe
\texttt{docs/mediciones/sus/} en el repositorio}: la prueba SUS a\'un no se
ha ejecutado. No se reporta ninguna puntuaci\'on SUS en este documento.

% EVIDENCIA-ZAIDA-05
% Ruta esperada: figuras/zaida/05-sus-resultados.png
\evidencia{figuras/zaida/05-sus-resultados.png}%
{Resultados agregados del instrumento SUS (media, desviaci\'on t\'ipica, IC 95\,\%), cuando est\'en disponibles.}%
{fig:zaida-sus-resultados}

\section{Accesibilidad autom\'atica (Lighthouse)}

El repositorio registra la intenci\'on de auditar el frontend con
Lighthouse~\cite{lighthouse-docs} CI (rama hist\'orica
\texttt{zaida/frontend-lighthouse}), pero
\textbf{no existe actualmente un \texttt{lighthouserc.js} versionado ni una
carpeta \texttt{docs/mediciones/lighthouse/}} con resultados. No se reporta
ninguna puntuaci\'on de Performance, Accessibility, Best Practices ni SEO
en este documento por no existir todav\'ia esa medici\'on.

% EVIDENCIA-ZAIDA-04
% Ruta esperada: figuras/zaida/04-lighthouse.png
\evidencia{figuras/zaida/04-lighthouse.png}%
{Reporte de Lighthouse CI del frontend de BIOPET, cuando est\'e disponible.}%
{fig:zaida-lighthouse}

% EVIDENCIA-ZAIDA-06
% Ruta esperada: figuras/zaida/06-accesibilidad.png
\evidencia{figuras/zaida/06-accesibilidad.png}%
{Detalle de accesibilidad autom\'atica (categor\'ia Accessibility de Lighthouse), cuando est\'e disponible.}%
{fig:zaida-accesibilidad}

\section{Publicabilidad y documentaci\'on funcional}

El SRS, las historias de usuario y los casos de uso est\'an versionados en
formato Markdown (migrados desde PDF en la consolidaci\'on m\'as reciente),
lo que permite su revisi\'on por diferencias (\emph{diff}) igual que el
resto del c\'odigo, en lugar de binarios opacos.

\section{Limitaciones reales de este bloque}

\begin{itemize}[nosep]
  \item Las mediciones de usabilidad (SUS) y accesibilidad autom\'atica
    (Lighthouse) est\'an dise\~nadas metodol\'ogicamente pero no ejecutadas
    a la fecha de este informe.
  \item La verificaci\'on de accesibilidad reportada aqu\'i es de
    inspecci\'on de c\'odigo (atributos ARIA presentes), no una auditor\'ia
    automatizada ni una prueba con usuarios reales con discapacidad.
  \item No se dispone de un reporte de compatibilidad entre navegadores
    distinto del entorno de desarrollo usado por el equipo.
\end{itemize}
